\parthead{I}{The airport that challenged intuition}\label{part:1}
From the ground, a large aircraft turning above a city looks like a contest between scale and space. At Kai Tak the setting made that impression unusually vivid: built-up Kowloon on the approach side, rising terrain beyond, and a runway extending into Kowloon Bay. The airport was not simply a long strip of pavement. Its usefulness depended on how aircraft could reach it.

\fig{01_historical_airport}{.79}{The airport and its immediate urban setting in Lands Department frame A36143. The annotations identify visible features; no flight track is overlaid. Catalogue date: 2 November 1993; the handwritten frame label appears to read 3 November, an unresolved source-date discrepancy. Original scan orientation, not north-up; no metric scale is inferred from this perspective photograph. Source: Lands Department, HKSAR Government; CSDI Portal and DATA.GOV.HK, open-data DAP-L0. Intellectual-property ownership retained by the Government; supplied as-is, without endorsement.}{fig:airport}

A runway-aligned instrument approach would have pointed aircraft towards a different surrounding landscape. Runway 13 instead used an offset Instrument Guidance System (IGS), followed by visual manoeuvring. The Civil Aviation Department dates the IGS introduction to 1974; it improved runway use in adverse weather but did not remove the final visual requirement \citep{cadhistory,cad}.

The question is not whether a 747 can turn. It plainly can. The question is whether the turn can finish at the right place, height, speed and orientation, with enough runway left for the rest of the landing.
\par\medskip
\subsection*{What the Checkerboard approach actually asked of a pilot}
The IGS supplied lateral guidance along a localiser and vertical guidance on a $3.1^\circ$ glidepath. In the 1993 report, its localiser centreline is $088^\circ$ magnetic; the runway direction is $135^\circ$ magnetic. Both bearings use the same reference, so their difference is $47^\circ$. That is an angular change in direction on a plan, not a command to bank the wings by 47 degrees.

\fig{02_approach_geometry}{1}{Published directions and the instrument-to-visual sequence. Bearings are clockwise from magnetic north. The rays share an origin only to compare angles: their meeting point is not a surveyed turn location, and the diagram is not a navigational chart. Source: CAD Report 1/95, sections 1.8.1 and 1.10, Appendix 1. Ground track and air heading differ in wind.}{fig:geometry}

The instrument segment ended at the middle marker, reported as 1.7 nautical miles from touchdown. If visual flight had not been achieved there, the report specifies a missed approach \citep[§1.8.1]{cad}. The famous final turn therefore was not something the IGS could guide to completion in zero visibility. Checkerboard Hill was a visual reference in Kowloon, not a landmark on Victoria Peak across the harbour. The CAD's retrospective account describes the visual right turn over Kowloon City \citep{cadlink}.

The approach distributed different jobs among navigation equipment, a published procedure, visual references and the crew. Terrain and buildings shaped the available geometry. Wind altered the relation between where the nose pointed and where the aircraft travelled. A successful alignment still had to deliver an acceptable touchdown.

The 1993 case report gives 3,331.5 m of runway length, 61 m width and 2,786 m of landing distance available (LDA) \citep[§1.10]{cad}. The last number governs the threshold-relative stopping calculation here. A general airport history lists a 3,390 m extension-era length \citep{cadhistory}; these source quantities are not silently substituted for one another. Likewise, the airport that closed in July 1998 must not be confused with the replacement at Chek Lap Kok \citep{cadlink}.

\insight{The mathematical question}{Can one continuous motion satisfy the demands of turning, descent, alignment, touchdown and stopping together? A solution must fit all of them, not merely look plausible from one viewpoint.}
\par\medskip
